% !TeX root = ../main.tex
%------------------------------------------------------------------------------
%  In-Situ Stress
%  Sources: Amadei & Stephansson (1997); Sheorey (1994); Hoek & Brown (1980);
%  ISRM Suggested Methods for stress determination.
%------------------------------------------------------------------------------
\chapter{In-Situ Stress}
\label{ch:in-situ-stress}

In-situ stress supplies the demand against which strength, classification and
support capacity are assessed. Neither side of that comparison carries meaning
without the other, yet stress remains the input most often assumed rather than
determined. The assumption ordinarily made, that vertical stress equals the
overburden weight and horizontal stress follows from an assumed ratio, is
adequate at moderate depth in simple settings and can be seriously in error
elsewhere. This chapter examines what determines the field, how it is
measured, how far the measurements can be relied upon, and how a design is
justified when the result remains uncertain.

\section{The stress field and its estimation}
\label{sec:stress-field}

\subsection{Origin and components}
\label{sec:stress-components}

The stress at a point in the ground is a second-order tensor with six
independent components, and a complete determination requires all six. Four
physical contributions combine to produce it.

Gravitational loading supplies the vertical component directly and, through
lateral restraint, a horizontal component. For a laterally confined elastic
half-space with no tectonic contribution, the horizontal-to-vertical ratio
follows from the Poisson effect,
\[
k = \frac{\nu}{1-\nu},
\]
which for $\nu$ between 0.2 and 0.3 gives $k$ between 0.25 and 0.43. Measured
ratios in the upper few hundred metres of the crust routinely exceed unity, so
this expression is better regarded as a lower bound than as an estimate.

Tectonic stress transmitted from plate boundaries through the lithosphere
supplies most of the excess. It is coherent over distances of hundreds of
kilometres, which is what makes regional compilations useful, and it is
largely independent of depth, so its relative importance diminishes as the
gravitational component grows.

Topography adds no stress of its own but redistributes the gravitational and
tectonic components. Stresses concentrate beneath valley
floors and relax beneath ridges, and the perturbation extends to a depth
comparable with the width of the topographic feature. This matters for
portals, for shallow alignments in mountainous terrain, and for any excavation
beneath a steep slope, where the principal directions may also rotate away
from the vertical.

Residual stress locked in at grain and block scale by cooling, diagenesis and
past deformation accounts for a smaller part of the total. Its practical
signature is in observations that otherwise appear anomalous: core discing,
time-dependent deformation of specimens after recovery, and measurement
scatter beyond instrumental precision.

Erosional unloading deserves particular attention, because it explains the
most consequential departure from the naive estimate. As overburden is
removed, the vertical stress falls with it, while horizontal strain remains
restrained and the horizontal stress falls more slowly. The result near
surface is a horizontal stress that can be several times the vertical
(Figure~\ref{fig:erosion-unloading}), the
reverse of what elastic confinement alone would predict. Shallow rock
excavations in formerly deeply buried terrain therefore behave as though they
were considerably deeper, and this is the usual explanation for spalling
encountered where none was anticipated.

\begin{figure}[htbp]
\centering
\datafigure[0.62\textwidth]{ch08/erosion-unloading}
\caption{Stress ratio after erosion of a thickness $\Delta z$ of overburden from ground with an initial $k = 1$, for a laterally restrained elastic layer with $\nu = 0.25$: the vertical stress falls by $\gamma\,\Delta z$ and the horizontal by only $\nu/(1-\nu)$ of that, so $k = 1 + (1 - \nu/(1-\nu))\,\Delta z/z$. Near surface $k$ reaches several times unity.}
\label{fig:erosion-unloading}
\end{figure}

\subsection{Magnitude with depth}
\label{sec:k-ratio}

The vertical stress is normally taken as $\sv = \gamma z$, with $\gamma$ the
unit weight of the overburden and $z$ the depth, and it is the most
reliable component of the field. The approximation fails where topography is
steep, beneath a valley or an overhang, and where a stiff unit bridges over a
soft one.

The horizontal stress is expressed through $k = \sh / \sv$. Compilations of
measurements show a characteristic pattern: $k$ is high and widely scattered
near surface, commonly between 1 and 3 and occasionally higher, and both the
value and its scatter fall with depth as the gravitational contribution comes
to dominate.
\citet{hoekbrown1980underground} bounded the observed data by
\[
\frac{100}{z} + 0.3 \;<\; k \;<\; \frac{1500}{z} + 0.5,
\]
with $z$ in metres. The envelope is wide, and deliberately so: at 100\,m it
admits $k$ between 1.3 and 15.5, and at 1000\,m between 0.4 and 2.0. Its value
lies in establishing a plausible range rather than an estimate.

\citet{sheorey1994theory} derived a physically based alternative from
a thermo-elastic model of the crust,
\[
k = 0.25 + 7 E_h \left( 0.001 + \frac{1}{z} \right),
\]
in which $E_h$ is the horizontal deformation modulus in \si{\giga\pascal} and
$z$ the depth in metres. The expression reproduces the observed hyperbolic
decay without empirical fitting and links $k$ to a measurable rock property,
which is its principal advantage (Figure~\ref{fig:k-depth-envelope}). It is
dimensionally inhomogeneous, as the
statement of units indicates, and it applies to the modulus of the rock mass
rather than of an intact specimen.

\begin{figure}[htbp]
\centering
\datafigure[\textwidth]{ch08/k-depth-envelope}
\caption{Stress ratio (left) and horizontal stress (right) with depth. Shaded: the envelope of \citet{hoekbrown1980underground}; broken curves: the relation of \citet{sheorey1994theory} for three moduli $E_h$; dotted: the confined elastic value for $\nu = 0.25$. Stresses for $\gamma = 27$~kN/m\textsuperscript{3}. At 1000~m the envelope still spans $k$ from 0.4 to 2.0.}
\label{fig:k-depth-envelope}
\end{figure}

Both relations should be treated as first estimates. Where stress governs the
design, in deep caverns in brittle rock or wherever spalling is anticipated,
direct determination is warranted.

\subsection{Measurement methods}
\label{sec:measurement-methods}

Three families of technique are in routine use~\citep{amadei1997stress}, and
they differ in what they
determine, not merely in how they determine it.

Overcoring can recover the complete stress tensor, in favourable cases from
a single measurement. A
small-diameter hole is instrumented, then overcored so that the surrounding
annulus of rock relaxes (Figure~\ref{fig:overcoring-schematic}), and the
recovered strains are converted to stresses
using elastic properties determined separately on the same core. The CSIRO
hollow inclusion cell, the USBM deformation gauge and the doorstopper are the
established variants; the hollow inclusion cell recovers all six components
from a single installation, whereas the deformation gauge and the doorstopper
each resolve only the stress in the plane normal to the hole and so require
measurements in three non-parallel holes. The method needs competent, reasonably intact rock, since
core that discs or fractures on relief yields no usable strain record, and
this excludes precisely the highly stressed ground in which the measurement is
most wanted. Depth is limited by the requirement for access to the collar of
the measurement hole.

\begin{figure}[htbp]
\centering
\datafigure[\textwidth]{ch08/overcoring-schematic}
\caption{Overcoring (schematic). Left: a strain cell bonded in a pilot hole of about 38~mm diameter is overcored at about 150~mm, isolating a core around it. Right: strains at the cell gauges as the bit advances past the gauge plane, $x$ being the bit position and $D$ the overcore diameter; the recovered strains are the stable values once the bit is about one diameter beyond.}
\label{fig:overcoring-schematic}
\end{figure}

Hydraulic fracturing pressurises an isolated borehole interval to the
breakdown pressure, at which a tensile fracture initiates, and then records
the shut-in pressure, at which the fracture closes, and the lower pressure at
which it reopens on later cycles (Figure~\ref{fig:hydrofrac-record}). The
shut-in pressure equals the stress normal
to the fracture, which for a vertical fracture is the minimum horizontal
stress, and this is the most reliable quantity the method produces. The
maximum horizontal stress is inferred either from the breakdown pressure,
through a relation that requires the tensile strength and the pore pressure,
both of which carry significant uncertainty, or from the reopening pressure,
which removes the tensile strength from the relation at the cost of
sensitivity to how the reopening point is picked on the pressure record. The
method assumes that one principal
stress is vertical, an assumption that is reasonable in flat-lying ground away
from major structure and unreliable elsewhere. Its advantages are depth
capability and tolerance of poor rock. Hydraulic testing of pre-existing
fractures relaxes the initiation assumption at the cost of requiring several
fractures of suitably varied orientation.

\begin{figure}[htbp]
\centering
\datafigure[\textwidth]{ch08/hydrofrac-record}
\caption{Pressure record of a hydraulic fracturing test (schematic): breakdown $P_b$ on the first cycle, reopening $P_r$ where a later cycle leaves the first-cycle slope (broken line), and shut-in $P_s$ after pumping stops. Drawn for $\sh = 10$, $\sigma_H = 16$, $T = 5$ and $p_0 = 1$~MPa, with $P_s = \sh$, $P_b = 3\sh - \sigma_H + T - p_0$ and $P_r = P_b - T$~\citep{amadei1997stress}.}
\label{fig:hydrofrac-record}
\end{figure}

Relief methods measure stress at an existing excavation surface. A flat jack
is installed in a slot cut in the wall, the closure across the slot recorded,
and the jack then repressurised until the original geometry is restored, at
which point the jack pressure equals the stress normal to the slot. The
measurement is direct and mechanically simple, and it samples the disturbed
zone rather than the far field. That is what is wanted for assessing an
existing lining or the present stability of a wall, and not what is wanted
for predicting the response of a new excavation, for which the far field is
the input. Each slot yields only the stress component normal to it, so slots
in several orientations are needed for even a partial tensor.

\section{Indirect evidence and back-analysis}
\label{sec:indirect-evidence}

Where direct measurement is unavailable, several classes of observation
constrain the field, and their combined use is often more informative than any
single determination.

Borehole breakout is the most valuable. Where the tangential stress at the
borehole wall exceeds the rock strength, the wall spalls in the direction of
the minimum horizontal stress, elongating the hole in a plane perpendicular to
the maximum horizontal stress. Orientation is recovered robustly by acoustic
or optical televiewer logging, and the breakout width constrains magnitude
through the Kirsch stress concentration~\citep{kirsch1898} and an assumed
strength (Section~\ref{sec:kirsch-appendix}): for a given strength, a wider
breakout implies a larger difference between the two horizontal stresses
(Figure~\ref{fig:borehole-breakout}).
Drilling-induced tensile fractures form perpendicular to the breakouts and
provide an independent check on orientation
(Figure~\ref{fig:r3d-breakout-block}). Both can be obtained from
geotechnical boreholes at modest additional cost, and both remain
systematically under-used.

\begin{figure}[htbp]
\centering
\datafigure[\textwidth]{ch08/borehole-breakout}
\caption{Borehole breakout. Left: breakouts form in the direction of $\sh$ and drilling-induced tensile fractures in the direction of $\sigma_H$ (schematic). Right: Kirsch stress at the wall, $\sigma_\theta = \sigma_H + \sh - 2(\sigma_H - \sh)\cos 2\theta$, for $\sh = 20$~MPa and no fluid pressure; for a wall strength $C_0 = 80$~MPa the breakout widens from $34^\circ$ to $73^\circ$ as $\sigma_H$ rises from 35 to 45~MPa.}
\label{fig:borehole-breakout}
\end{figure}

\begin{figure}[htbp]
\centering
\renderfigure{breakout_block}{%
  \rlab[west]{0.53,0.72}{0.72,0.90}{Breakout}
  \rlab[east]{0.435,0.60}{0.28,0.86}{Tensile fracture}
}
\caption{A vertical borehole in a rock cube, a quarter cut away above a
horizontal slice. Breakouts (carmine) are spalled notches centred on the $\sh$
azimuth, $73^\circ$ wide for $\sigma_H = 45$, $\sh = 20$ and $C_0 = 80$~MPa
(Figure~\ref{fig:borehole-breakout}); drilling-induced tensile fractures
(malachite) are thin traces at the $\sigma_H$ azimuth. Notch and fracture
depths, 0.45 and 0.35 of the radius, are schematic.}
\label{fig:r3d-breakout-block}
\end{figure}

Core discing indicates high stress relative to strength. It is a qualitative
indicator, but its onset along a borehole marks the depth below which
stress-driven behaviour should be anticipated, and disc thickness and
curvature carry some information about magnitude. Focal mechanisms from local
seismicity and slip data from mapped faults constrain the regional field over
much larger volumes than any borehole measurement.

Back-analysis from observed behaviour becomes available once construction
begins, and it frequently yields the best estimate a project ever obtains.
Measured convergence, the pattern and depth of overbreak, and the position of
spalling around the profile all depend on the stress field, and a model
calibrated to reproduce them constrains it. The limitation is non-uniqueness:
stress, strength and modulus trade against one another, and several
combinations reproduce the same observation. Back-analysis is therefore
reliable when it refines a field already bracketed by measurement, and
unreliable when asked to determine stress and strength simultaneously from
displacement alone. Calibrating against several independent measured
quantities rather than one substantially reduces the ambiguity.

\section{Designing under stress uncertainty}
\label{sec:stress-uncertainty}

Stress measurements scatter~\citep{amadei1997stress}. Repeat determinations at
one site commonly differ by 20 to 50\,per cent in magnitude and by tens of
degrees in orientation, and
the scatter is largely real rather than instrumental. It reflects genuine
variability at the scale of the measurement in a quantity that is sampled at
points and applied over kilometres, so additional measurements reduce the
standard error of the mean without reducing the variability itself.

The design response is not to seek a single best estimate but to construct an
envelope and test against it. An upper-bound case with high horizontal stress
governs spalling, stress-induced failure in the sidewalls and support demand
in the crown. A lower-bound case with low horizontal stress governs tensile
zones in the crown, structurally controlled failure and, in shallow works,
uplift (Figure~\ref{fig:stress-bracket-kirsch}). Because different bounds
govern different mechanisms, neither can be
described as conservative in general, and a design justified only at the mean
has been checked against no mechanism at all.

\begin{figure}[htbp]
\centering
\datafigure[0.62\textwidth]{ch08/stress-bracket-kirsch}
\caption{Kirsch boundary stress at the crown, $(3k-1)\sv$, and sidewall, $(3-k)\sv$, of a circular opening (Section~\ref{sec:kirsch-appendix}). Shaded: the range of $k$ admitted by the depth envelope at 1000~m, across which the crown stress varies 25-fold; below $k = 1/3$ the crown is in tension (tinted band).}
\label{fig:stress-bracket-kirsch}
\end{figure}

Orientation carries its own consequence, and it is the one variable a designer
can sometimes exploit. The relationship between the excavation axis and the
principal stress directions determines the boundary stress distribution around
the profile, and in a strongly anisotropic field the difference between a
favourable and an unfavourable orientation can be the difference between
routine support and a persistent spalling problem. Where the alignment has
freedom this becomes a design decision, treated in
Chapter~\ref{ch:openings}; where it does not, it is a constraint that should be
recognised early enough to influence the choice of excavation method.

A defensible first pass proceeds as follows. Take $\sv = \gamma z$ and check it
against the topography. Bracket $k$ using the depth envelope and, where rock
properties allow, the thermo-elastic estimate. Take orientation from the
regional compilation and check it against any breakout data available in the
project boreholes. State the resulting field as a range rather than as a
value. Then ask whether the design is sensitive to that range. If the support
class, the predicted extent of the spalling-prone zone and the crown tensile
zone are all unchanged across it, the estimate is adequate and no measurement
programme is required; if it is not, that sensitivity is itself the
justification for one. The same test governs investigation targeting in
Chapter~\ref{ch:ground-modelling}, and it is the appropriate basis for
deciding how much stress determination a project needs.

\chapterreferences
